% Options for packages loaded elsewhere
\PassOptionsToPackage{unicode}{hyperref}
\PassOptionsToPackage{hyphens}{url}
%
\documentclass[
]{article}
\usepackage{amsmath,amssymb}
\usepackage{iftex}
\ifPDFTeX
  \usepackage[T1]{fontenc}
  \usepackage[utf8]{inputenc}
  \usepackage{textcomp} % provide euro and other symbols
\else % if luatex or xetex
  \usepackage{unicode-math} % this also loads fontspec
  \defaultfontfeatures{Scale=MatchLowercase}
  \defaultfontfeatures[\rmfamily]{Ligatures=TeX,Scale=1}
\fi
\usepackage{lmodern}
\ifPDFTeX\else
  % xetex/luatex font selection
\fi
% Use upquote if available, for straight quotes in verbatim environments
\IfFileExists{upquote.sty}{\usepackage{upquote}}{}
\IfFileExists{microtype.sty}{% use microtype if available
  \usepackage[]{microtype}
  \UseMicrotypeSet[protrusion]{basicmath} % disable protrusion for tt fonts
}{}
\makeatletter
\@ifundefined{KOMAClassName}{% if non-KOMA class
  \IfFileExists{parskip.sty}{%
    \usepackage{parskip}
  }{% else
    \setlength{\parindent}{0pt}
    \setlength{\parskip}{6pt plus 2pt minus 1pt}}
}{% if KOMA class
  \KOMAoptions{parskip=half}}
\makeatother
\usepackage{xcolor}
\usepackage[margin=1in]{geometry}
\usepackage{color}
\usepackage{fancyvrb}
\newcommand{\VerbBar}{|}
\newcommand{\VERB}{\Verb[commandchars=\\\{\}]}
\DefineVerbatimEnvironment{Highlighting}{Verbatim}{commandchars=\\\{\}}
% Add ',fontsize=\small' for more characters per line
\usepackage{framed}
\definecolor{shadecolor}{RGB}{248,248,248}
\newenvironment{Shaded}{\begin{snugshade}}{\end{snugshade}}
\newcommand{\AlertTok}[1]{\textcolor[rgb]{0.94,0.16,0.16}{#1}}
\newcommand{\AnnotationTok}[1]{\textcolor[rgb]{0.56,0.35,0.01}{\textbf{\textit{#1}}}}
\newcommand{\AttributeTok}[1]{\textcolor[rgb]{0.13,0.29,0.53}{#1}}
\newcommand{\BaseNTok}[1]{\textcolor[rgb]{0.00,0.00,0.81}{#1}}
\newcommand{\BuiltInTok}[1]{#1}
\newcommand{\CharTok}[1]{\textcolor[rgb]{0.31,0.60,0.02}{#1}}
\newcommand{\CommentTok}[1]{\textcolor[rgb]{0.56,0.35,0.01}{\textit{#1}}}
\newcommand{\CommentVarTok}[1]{\textcolor[rgb]{0.56,0.35,0.01}{\textbf{\textit{#1}}}}
\newcommand{\ConstantTok}[1]{\textcolor[rgb]{0.56,0.35,0.01}{#1}}
\newcommand{\ControlFlowTok}[1]{\textcolor[rgb]{0.13,0.29,0.53}{\textbf{#1}}}
\newcommand{\DataTypeTok}[1]{\textcolor[rgb]{0.13,0.29,0.53}{#1}}
\newcommand{\DecValTok}[1]{\textcolor[rgb]{0.00,0.00,0.81}{#1}}
\newcommand{\DocumentationTok}[1]{\textcolor[rgb]{0.56,0.35,0.01}{\textbf{\textit{#1}}}}
\newcommand{\ErrorTok}[1]{\textcolor[rgb]{0.64,0.00,0.00}{\textbf{#1}}}
\newcommand{\ExtensionTok}[1]{#1}
\newcommand{\FloatTok}[1]{\textcolor[rgb]{0.00,0.00,0.81}{#1}}
\newcommand{\FunctionTok}[1]{\textcolor[rgb]{0.13,0.29,0.53}{\textbf{#1}}}
\newcommand{\ImportTok}[1]{#1}
\newcommand{\InformationTok}[1]{\textcolor[rgb]{0.56,0.35,0.01}{\textbf{\textit{#1}}}}
\newcommand{\KeywordTok}[1]{\textcolor[rgb]{0.13,0.29,0.53}{\textbf{#1}}}
\newcommand{\NormalTok}[1]{#1}
\newcommand{\OperatorTok}[1]{\textcolor[rgb]{0.81,0.36,0.00}{\textbf{#1}}}
\newcommand{\OtherTok}[1]{\textcolor[rgb]{0.56,0.35,0.01}{#1}}
\newcommand{\PreprocessorTok}[1]{\textcolor[rgb]{0.56,0.35,0.01}{\textit{#1}}}
\newcommand{\RegionMarkerTok}[1]{#1}
\newcommand{\SpecialCharTok}[1]{\textcolor[rgb]{0.81,0.36,0.00}{\textbf{#1}}}
\newcommand{\SpecialStringTok}[1]{\textcolor[rgb]{0.31,0.60,0.02}{#1}}
\newcommand{\StringTok}[1]{\textcolor[rgb]{0.31,0.60,0.02}{#1}}
\newcommand{\VariableTok}[1]{\textcolor[rgb]{0.00,0.00,0.00}{#1}}
\newcommand{\VerbatimStringTok}[1]{\textcolor[rgb]{0.31,0.60,0.02}{#1}}
\newcommand{\WarningTok}[1]{\textcolor[rgb]{0.56,0.35,0.01}{\textbf{\textit{#1}}}}
\usepackage{graphicx}
\makeatletter
\def\maxwidth{\ifdim\Gin@nat@width>\linewidth\linewidth\else\Gin@nat@width\fi}
\def\maxheight{\ifdim\Gin@nat@height>\textheight\textheight\else\Gin@nat@height\fi}
\makeatother
% Scale images if necessary, so that they will not overflow the page
% margins by default, and it is still possible to overwrite the defaults
% using explicit options in \includegraphics[width, height, ...]{}
\setkeys{Gin}{width=\maxwidth,height=\maxheight,keepaspectratio}
% Set default figure placement to htbp
\makeatletter
\def\fps@figure{htbp}
\makeatother
\setlength{\emergencystretch}{3em} % prevent overfull lines
\providecommand{\tightlist}{%
  \setlength{\itemsep}{0pt}\setlength{\parskip}{0pt}}
\setcounter{secnumdepth}{-\maxdimen} % remove section numbering
\ifLuaTeX
  \usepackage{selnolig}  % disable illegal ligatures
\fi
\usepackage{bookmark}
\IfFileExists{xurl.sty}{\usepackage{xurl}}{} % add URL line breaks if available
\urlstyle{same}
\hypersetup{
  pdftitle={tø uge 2},
  hidelinks,
  pdfcreator={LaTeX via pandoc}}

\title{tø uge 2}
\author{}
\date{\vspace{-2.5em}2025-01-30}

\begin{document}
\maketitle

Compute the mean x and median m of the eight numbers
\(1, 2, 4, 5, 6, 8, 11\) 15. Let \[f (x) = \sqrt{x}\].\newline

Apply this transformation to the data, and then com- pute the mean \(x̃\)
and the median \(m̃\) of the transformed data

question a: \[\text{ Is}\,\, f (x) = x̃?\]\newline question b:
\[\text{Is}  f ( m ) = m̃?\]

\begin{Shaded}
\begin{Highlighting}[]
\CommentTok{\#we write up the data }
\NormalTok{x}\OtherTok{=}\FunctionTok{c}\NormalTok{(}\DecValTok{1}\NormalTok{,}\DecValTok{2}\NormalTok{,}\DecValTok{4}\NormalTok{,}\DecValTok{5}\NormalTok{,}\DecValTok{6}\NormalTok{,}\DecValTok{8}\NormalTok{,}\DecValTok{11}\NormalTok{)}
\CommentTok{\#calculate the transformed data}
\NormalTok{fx}\OtherTok{=}\FunctionTok{sqrt}\NormalTok{(x)}
\CommentTok{\# find the mean for the data}
\NormalTok{xb}\OtherTok{=}\FunctionTok{mean}\NormalTok{(x)}
\CommentTok{\#then the mean of the transformation }
\FunctionTok{mean}\NormalTok{(fx)}
\end{Highlighting}
\end{Shaded}

\begin{verbatim}
## [1] 2.177832
\end{verbatim}

\begin{Shaded}
\begin{Highlighting}[]
\CommentTok{\#then f of the mean }
\FunctionTok{sqrt}\NormalTok{(xb)}
\end{Highlighting}
\end{Shaded}

\begin{verbatim}
## [1] 2.299068
\end{verbatim}

\begin{Shaded}
\begin{Highlighting}[]
\CommentTok{\#now we do the exactly the same but for the median }
\NormalTok{m}\OtherTok{=}\FunctionTok{median}\NormalTok{(x)}
\FunctionTok{median}\NormalTok{(fx)}
\end{Highlighting}
\end{Shaded}

\begin{verbatim}
## [1] 2.236068
\end{verbatim}

\begin{Shaded}
\begin{Highlighting}[]
\FunctionTok{sqrt}\NormalTok{(m)}
\end{Highlighting}
\end{Shaded}

\begin{verbatim}
## [1] 2.236068
\end{verbatim}

the mean is not equal since f is concave then\\
\[f(\dfrac{1}{n}*\sum_{i=1}^nx_i)\leq \dfrac{1}{n} \sum_{i=1}^nf(x_i)\]
the median in this case is equal but this is only true in general when f
in linear or n is uneven and f is montone

Let x and m denote the mean and median, respectively, of
\[x_1 < x_2 < · · · < x_n\] Let f be a real-valued function.

\begin{enumerate}
\def\labelenumi{\alph{enumi})}
\item
  Is f(x) the mean off \[f(x_1),f(x_2),...,f(x_n)\]? no simply look at
  the last exercise
\item
  Is f(m) the median off f(x1),f(x2),\ldots,f(xn)? try
  \(x=(x_1,x_2)=(1,2)\) \(f=x^2\) not equal
\item
  Are there any conditions that would ensure that f (x) is the mean of
  the transformed data? liniarity of f or no variance in x
\item
  Are there any conditions that would ensure that f (m) is the median of
  the transformed data?
\end{enumerate}

liniarity of f or no variance in x or f order preserving and n uneven

\begin{Shaded}
\begin{Highlighting}[]
\NormalTok{s}\OtherTok{=}\FunctionTok{table}\NormalTok{(FlightDelays}\SpecialCharTok{$}\NormalTok{DepartTime)}
\FunctionTok{barplot}\NormalTok{(s)}
\end{Highlighting}
\end{Shaded}

\includegraphics{tø-uge-2_files/figure-latex/unnamed-chunk-2-1.pdf}

\begin{Shaded}
\begin{Highlighting}[]
\NormalTok{tab }\OtherTok{\textless{}{-}} \FunctionTok{table}\NormalTok{(FlightDelays}\SpecialCharTok{$}\NormalTok{Day,FlightDelays}\SpecialCharTok{$}\NormalTok{Delayed30)}
\FunctionTok{prop.table}\NormalTok{(tab,}\DecValTok{1}\NormalTok{)}
\end{Highlighting}
\end{Shaded}

\begin{verbatim}
##      
##               No        Yes
##   Fri 0.77394035 0.22605965
##   Mon 0.90317460 0.09682540
##   Sat 0.89624724 0.10375276
##   Sun 0.92014519 0.07985481
##   Thu 0.76678445 0.23321555
##   Tue 0.85191083 0.14808917
##   Wed 0.86524823 0.13475177
\end{verbatim}

\begin{Shaded}
\begin{Highlighting}[]
\FunctionTok{boxplot}\NormalTok{(FlightDelays}\SpecialCharTok{$}\NormalTok{FlightLength}\SpecialCharTok{\textasciitilde{}}\NormalTok{FlightDelays}\SpecialCharTok{$}\NormalTok{Delayed30)}
\end{Highlighting}
\end{Shaded}

\includegraphics{tø-uge-2_files/figure-latex/unnamed-chunk-4-1.pdf}

opgave 2.4,d vi ser at hvis flyet er forsiket er der en støøre chance
for at det er en medium længde flyvetur men også mange flere outliers,
jeg ville kke drage nogle konklutioner

\begin{Shaded}
\begin{Highlighting}[]
\FunctionTok{par}\NormalTok{(}\AttributeTok{mfrow=}\FunctionTok{c}\NormalTok{(}\DecValTok{1}\NormalTok{,}\DecValTok{2}\NormalTok{))}
\NormalTok{tab }\OtherTok{\textless{}{-}} \FunctionTok{table}\NormalTok{(Recidivism}\SpecialCharTok{$}\NormalTok{Recid)}
\FunctionTok{barplot}\NormalTok{(tab)}
\FunctionTok{prop.table}\NormalTok{(tab)}
\end{Highlighting}
\end{Shaded}

\begin{verbatim}
## 
##        No       Yes 
## 0.6835859 0.3164141
\end{verbatim}

\begin{Shaded}
\begin{Highlighting}[]
\NormalTok{tab2 }\OtherTok{\textless{}{-}}\FunctionTok{table}\NormalTok{(Recidivism}\SpecialCharTok{$}\NormalTok{Age25,Recidivism}\SpecialCharTok{$}\NormalTok{Recid)}
\FunctionTok{prop.table}\NormalTok{(tab2,}\DecValTok{1}\NormalTok{)}
\end{Highlighting}
\end{Shaded}

\begin{verbatim}
##           
##                   No       Yes
##   Over 25  0.6942333 0.3057667
##   Under 25 0.6350341 0.3649659
\end{verbatim}

\includegraphics{tø-uge-2_files/figure-latex/unnamed-chunk-5-1.pdf}

wee seee that about 30,6\% of offenders over 25 go back to prison while
about 36.5\% of under 25 offenders go back

\begin{Shaded}
\begin{Highlighting}[]
\NormalTok{t}\OtherTok{=}\NormalTok{Recidivism}\SpecialCharTok{$}\NormalTok{Days}
\NormalTok{s}\OtherTok{=}\NormalTok{Recidivism}\SpecialCharTok{$}\NormalTok{Type}
\FunctionTok{boxplot}\NormalTok{(t}\SpecialCharTok{\textasciitilde{}}\NormalTok{s)}
\end{Highlighting}
\end{Shaded}

\includegraphics{tø-uge-2_files/figure-latex/unnamed-chunk-6-1.pdf}

the median reeviction time is lover for tech crimes . all quantiles are
lover for the tech crime type

\begin{Shaded}
\begin{Highlighting}[]
\FunctionTok{quantile}\NormalTok{(Recidivism}\SpecialCharTok{$}\NormalTok{Days,}\AttributeTok{na.rm=}\ConstantTok{TRUE}\NormalTok{)}
\end{Highlighting}
\end{Shaded}

\begin{verbatim}
##   0%  25%  50%  75% 100% 
##    0  241  418  687 1095
\end{verbatim}

\begin{Shaded}
\begin{Highlighting}[]
\FunctionTok{par}\NormalTok{(}\AttributeTok{mfrow=}\FunctionTok{c}\NormalTok{(}\DecValTok{1}\NormalTok{,}\DecValTok{2}\NormalTok{))}
\NormalTok{over25 }\OtherTok{\textless{}{-}}\FunctionTok{subset}\NormalTok{(Recidivism, Age25}\SpecialCharTok{==} \StringTok{"Over 25"}\NormalTok{)}
\NormalTok{under25 }\OtherTok{\textless{}{-}} \FunctionTok{subset}\NormalTok{(Recidivism,Age25}\SpecialCharTok{==} \StringTok{"Under 25"}\NormalTok{)}

\FunctionTok{plot.ecdf}\NormalTok{(over25}\SpecialCharTok{$}\NormalTok{Days)}
\FunctionTok{plot.ecdf}\NormalTok{(under25}\SpecialCharTok{$}\NormalTok{Days)}
\end{Highlighting}
\end{Shaded}

\includegraphics{tø-uge-2_files/figure-latex/unnamed-chunk-8-1.pdf}

opgave 1.26

\begin{Shaded}
\begin{Highlighting}[]
\FunctionTok{curve}\NormalTok{(}\FunctionTok{log}\NormalTok{(x),}\DecValTok{0}\NormalTok{,}\DecValTok{2}\NormalTok{,}\DecValTok{10000}\NormalTok{)}
\FunctionTok{abline}\NormalTok{(}\AttributeTok{h=}\DecValTok{0}\NormalTok{)}
\FunctionTok{abline}\NormalTok{(}\AttributeTok{v=}\DecValTok{1}\NormalTok{)}
\end{Highlighting}
\end{Shaded}

\includegraphics{tø-uge-2_files/figure-latex/unnamed-chunk-9-1.pdf}

\begin{Shaded}
\begin{Highlighting}[]
\FunctionTok{curve}\NormalTok{(x}\SpecialCharTok{*}\FunctionTok{sin}\NormalTok{(x}\SpecialCharTok{\^{}}\DecValTok{2}\NormalTok{),}\DecValTok{0}\NormalTok{,}\DecValTok{2}\SpecialCharTok{*}\NormalTok{pi,}\DecValTok{1000}\NormalTok{)}
\FunctionTok{curve}\NormalTok{(x}\SpecialCharTok{*}\FunctionTok{cos}\NormalTok{(x}\SpecialCharTok{\^{}}\DecValTok{2}\NormalTok{),}\DecValTok{0}\NormalTok{,}\DecValTok{2}\SpecialCharTok{*}\NormalTok{pi,}\DecValTok{1000}\NormalTok{,}\AttributeTok{add =} \ConstantTok{TRUE}\NormalTok{)}
\end{Highlighting}
\end{Shaded}

\includegraphics{tø-uge-2_files/figure-latex/unnamed-chunk-10-1.pdf}

opgave 1.27

\begin{Shaded}
\begin{Highlighting}[]
\NormalTok{data }\OtherTok{\textless{}{-}}\NormalTok{TXBirths2004}
\FunctionTok{hist}\NormalTok{(TXBirths2004}\SpecialCharTok{$}\NormalTok{Weight, }\AttributeTok{freq=} \ConstantTok{FALSE}\NormalTok{, }\AttributeTok{breaks=}\DecValTok{250}\SpecialCharTok{*}\DecValTok{0}\SpecialCharTok{:}\DecValTok{20}\NormalTok{)}


\FunctionTok{curve}\NormalTok{(}\FunctionTok{dnorm}\NormalTok{(x,}\AttributeTok{mean=}\FunctionTok{mean}\NormalTok{(TXBirths2004}\SpecialCharTok{$}\NormalTok{Weight),}\AttributeTok{sd=}\FunctionTok{sd}\NormalTok{(TXBirths2004}\SpecialCharTok{$}\NormalTok{Weight)),}\DecValTok{0}\NormalTok{,}\DecValTok{5000}\NormalTok{, }\AttributeTok{add=}\ConstantTok{TRUE}\NormalTok{)}
\end{Highlighting}
\end{Shaded}

\includegraphics{tø-uge-2_files/figure-latex/unnamed-chunk-11-1.pdf}

\begin{Shaded}
\begin{Highlighting}[]
\FunctionTok{hist}\NormalTok{(TXBirths2004}\SpecialCharTok{$}\NormalTok{Weight, }\AttributeTok{freq=} \ConstantTok{FALSE}\NormalTok{, }\AttributeTok{breaks=}\DecValTok{50}\SpecialCharTok{*}\DecValTok{0}\SpecialCharTok{:}\DecValTok{100}\NormalTok{)}
\FunctionTok{curve}\NormalTok{(}\FunctionTok{dnorm}\NormalTok{(x,}\AttributeTok{mean=}\FunctionTok{mean}\NormalTok{(TXBirths2004}\SpecialCharTok{$}\NormalTok{Weight),}\AttributeTok{sd=}\FunctionTok{sd}\NormalTok{(TXBirths2004}\SpecialCharTok{$}\NormalTok{Weight)),}\DecValTok{0}\NormalTok{,}\DecValTok{5000}\NormalTok{, }\AttributeTok{add=}\ConstantTok{TRUE}\NormalTok{)}
\end{Highlighting}
\end{Shaded}

\includegraphics{tø-uge-2_files/figure-latex/unnamed-chunk-12-1.pdf}

histogrammet bliver langt mindre overskueligt og mere flimret og derved
sværere at konkludere ud fra \}

oms opgave 1.25 part a)

\[\hat{F}_n(x)=\dfrac{1}{n}*\sum_{i=1}^n\mathbb{1}_{]-\infty,x]}(X_i)\]

vi noterer at:
\[\mathbb{1}_{]-\infty,x]}(X_i)=1\Leftrightarrow X_i\leq x\] og
\[\mathbb{1}_{]-\infty,x]}(X_i)=0\Leftrightarrow X_i> x\] da er denne
transformation bernulli med sandylingheds paramete
\[\mathbb{P}(X_i \leq x)=F(x)\] summen af i.i.d bernulli er binomial
\[n*\hat{F}_n(x)=\sum_{i=1}^n\mathbb{1}_{]-\infty,x]}(X_i)\] er da
binomial

del b)

\[\mathbb{E}[\hat{F}_n(x)]=\dfrac{1}{n}*\mathbb{E}[n*\hat{F}_n(x)]=F(x)\]

\[\mathbb{V}[\hat{F}_n(x)]=\dfrac{1}{n^2}*\mathbb{V}[n*\hat{F}_n(x)]=\dfrac{1}{n^2}*n*F(x)*[1-F(x)]\]

\[=\dfrac{1}{n}*F(x)*(1-F(x))\]

del c) vis

\[\lim_{n\to \infty}\mathbb{P}(\lvert \hat{F}_n(x)-F(x)\lvert \geq \delta)=0 \quad \forall  \delta>0\]

vi bemærker at
\[\hat{F}_n(x)=\dfrac{1}{n}*\sum_{i=1}^n\mathbb{1}_{]-\infty,x]}(X_i)\]
er en stikprøvemiddelværdi for varablen vi allerede har beskrevet som
bernulli da siger store tals lov
\[ \overline X \to^{\mathbb{P}} \mu \quad for\,\, n\to\infty \]

da \[\hat{F}_n(x)\] er en stkprøve middelværdi for bernulli må:
\[\hat{F}_n(x) \to^{\mathbb{P}}F(x)  \quad for\,\, n\to\infty \]

\end{document}
